%%% TUTAJ PODAJ LISTĘ SWOICH AKRONIMÓW
%%% Aby poprawnie wyświetliły wszyskie skróty
%%% należy kilka razy skompilować plik główny

\RequirePackage{acro}

%% Uważaj na formatowanie;
%% Nie zapomnij o przecinkach

\DeclareAcronym{dry}{
  short = DRY ,
  long  = Don't repeat yourself
}
\DeclareAcronym{kiss}{
  short = KISS ,
  long  = {Keep It Simple, Stupid}
}
\DeclareAcronym{sql}{
  short = SQL,
  long  = Structured Query Language
}
\DeclareAcronym{api}{
  short = API ,
  long  = Application Programming Interface,
  short-plural = APIs ,
  long-plural  = Application Programming Interfaces
}

%%%%%%%%%%%%%%%%%%%%%%%%%%%%%%%%%%%%%%%%%%%
%% NIE RUSZAJ TEGO CO JEST NIŻEJ JEŚLI  %%%
%% NIE ZNASZ ZBYT DOBRZE LATEX'A        %%%
%%%%%%%%%%%%%%%%%%%%%%%%%%%%%%%%%%%%%%%%%%%
\newcommand{\ShowAcronyms}[1]{
  \begin{center}
    {\Large\bfseries Lista skrótów \par}
  \end{center}
  \vspace{0.5em}
	
  \ifboolexpr{
      test {\ifstrequal{#1}{lof}} or
      test {\ifstrequal{#1}{toc}} or
      test {\ifstrequal{#1}{description}}
    }{%
      \printacronyms[sort=true, heading=none, template=#1]
    }{%
      \printacronyms[sort=true, heading=none, template=lof]
     }%
  \clearpage
}